\documentclass[t,20pt]{beamer}

% BiB corporate theme (local .sty files in this folder). The font theme loads
% fontspec, so this deck must be built with XeLaTeX (see .latexmkrc).
\usetheme{bibcd}

\usepackage[ngerman,english]{babel}
\usepackage{amsmath,amssymb}
\usepackage{booktabs}
\usepackage{tabularx}
\usetikzlibrary{arrows.meta,positioning,calc}

% Figures come straight from the analysis folder, so the slides always show
% the current versions.
\graphicspath{{../figures/}}

% Slightly larger list text than the template default, for readability in the room.
\setbeamerfont{itemize/enumerate body}{size=\normalsize}
\setbeamerfont{itemize/enumerate subbody}{size=\small}
\setbeamerfont{itemize/enumerate subsubbody}{size=\footnotesize}

\hypersetup{pdftitle={Divorce-induced health decline and the impact on mortality risk},
            pdfauthor={Kar Man Tan, Pavel Grigoriev}}

% Helpers
\newcommand{\hl}[1]{\textcolor{BiBViolett}{\textbf{#1}}}           % highlight
\newcommand{\ct}[1]{{\color{BiBBlau80}(#1)}}                        % in-text citation
\newcommand{\slidenote}[1]{\par\vfill{\footnotesize\color{gray}#1\par}} % note at slide bottom
\newcommand{\ci}[2]{{\footnotesize[$#1$, $#2$]}}                    % confidence interval
\newcommand{\Lam}{\Lambda}
% Smaller list text for dense slides. Beamer sets the list font itself, so a
% surrounding {\small ...} has no effect on itemize; use this at the start of a frame.
\newcommand{\compactlists}{%
  \setbeamerfont{itemize/enumerate body}{size=\small}%
  \setbeamerfont{itemize/enumerate subbody}{size=\footnotesize}}
\newcolumntype{Y}{>{\raggedright\arraybackslash}X}

\title{Divorce-induced health decline and the impact on mortality risk}
\subtitle{DIDI Workshop, Wiesbaden\\23 Nov 2026}
\author{Kar Man Tan and Pavel Grigoriev}

\begin{document}

% =====================================================================
\begin{bibframe}[frametype=titleen]
  \titlepage
\end{bibframe}

\begin{bibframe}[frametype=agenda]
  \frametitle{Contents}
  \begin{columns}
    \begin{column}{10.5cm}
    \end{column}
    \begin{column}{18.5cm}
      \tableofcontents[hideallsubsections]
    \end{column}
  \end{columns}
\end{bibframe}

% =====================================================================
\section{Motivation and introduction}

\begin{bibframe}
  \frametitle{Why divorce, health and mortality?}
  \begin{itemize}
    \item Divorce is consistently linked to higher mortality: mean hazard ratio \hl{1.30} across 104 studies \ct{Shor et al., 2012}
    \item The mechanisms are poorly identified
    \begin{itemize}
      \item Health is both a consequence and a cause of divorce \ct{Joung et al., 1998; Torvik et al., 2015}
      \item Who divorces is not random
    \end{itemize}
    \item Germany: very little evidence. Brockmann \& Klein (2004), using the SOEP survey: divorced men +51\% mortality, no significant effect for women
    \item Since 2011, DRV registers can be linked: income biographies, divorce dates (pension splitting), rehabilitation histories and deaths
    \item[$\rightarrow$] Project \hl{DiDi}: \textit{Divorce and Diagnosis} (BiB, Hertie School, Charité, DRV)
  \end{itemize}
\end{bibframe}

\begin{bibframe}
  \frametitle{Why musculoskeletal (MSK) rehabilitation?}
  \begin{itemize}
    \item MSK conditions are the largest source of lost work capacity in Germany
    \begin{itemize}
      \item 21\% of certified incapacity days in 2024, the largest diagnostic group \ct{BAuA, 2026}
      \item Orthopaedic conditions are the most common completed medical rehabilitation: 35.6\% (women), 35.5\% (men) \ct{DRV Reha-Atlas, 2024}
      \item In our data: 46.5\% of all completed rehabilitation spells
    \end{itemize}
    \item The DRV grants rehabilitation only once a condition is chronic and \hl{threatens the ability to work} \ct{DRV, 2005}
    \item Our outcome therefore measures crossing an \hl{institutional threshold}, not the onset of disease
    \begin{itemize}
      \item Does divorce push people with a latent vulnerability across that threshold?
    \end{itemize}
  \end{itemize}
\end{bibframe}

\begin{bibframe}
  \frametitle{Research questions and approach}
  \begin{block}{RQ1}
    Does divorce raise the risk of a first musculoskeletal rehabilitation?
  \end{block}
  \begin{block}{RQ2}
    Is part of the excess mortality after divorce due to divorce-induced musculoskeletal deterioration?
  \end{block}
  \vspace{0.3cm}
  \begin{itemize}
    \item Annual risk sets: first observed divorce in year $t_0$ vs.\ people not yet divorced in $t_0$
    \item Propensity score matching on up to three pre-divorce years
    \item Discrete-time hazards, cumulative incidence of rehabilitation, cumulative mortality
    \item Interventional mediation analysis via g-computation (RQ2)
  \end{itemize}
\end{bibframe}

% =====================================================================
\section{Literature review}

\subsection{Divorce and mortality}
\begin{bibframe}
  \frametitle{Literature: divorce and mortality}
  \begin{itemize}
    \item Meta-analysis: hazard ratio 1.30 for the divorced; men 1.37, women 1.22; the gender gap narrows with age \ct{Shor et al., 2012}
    \item Timing differs across countries
    \begin{itemize}
      \item Finland: men's excess mortality (accidents, violence, alcohol) is highest right after divorce \ct{Metsä-Simola \& Martikainen, 2013}
      \item Norway: men's excess grows with time since divorce; remarriage changes it \ct{Berntsen \& Kravdal, 2012}
    \end{itemize}
    \item Germany: selection into marriage; longer marriages go with lower mortality \ct{Brockmann \& Klein, 2004}
    \item Designs aimed at causality shrink the effect
    \begin{itemize}
      \item Exposure definition matters: continuously divorced HR 1.66 vs.\ ever divorced 0.98 \ct{Sbarra \& Nietert, 2009}
      \item Norway: HR 1.59 (men) and 1.47 (women) disappear within twin pairs \ct{Wiborg \& Gugushvili, 2025}
    \end{itemize}
  \end{itemize}
\end{bibframe}

\subsection{Divorce and health}
\begin{bibframe}
  \frametitle{Literature: divorce and health}
  \begin{itemize}
    \item Poor health often \hl{precedes} divorce
    \begin{itemize}
      \item Health complaints: 1.5$\times$, multiple chronic conditions: 2$\times$ more likely to divorce \ct{Joung et al., 1998}
      \item Health behaviours and distress predict divorce \ct{Torvik et al., 2015}; psychotropic use is elevated five years before divorce \ct{Metsä-Simola \& Martikainen, 2013}
    \end{itemize}
    \item Accounting for selection leaves small effects
    \begin{itemize}
      \item Marginal effect only for men under 45 \ct{Kohn \& Averett, 2014}
      \item After matching, only the previously depressed have depressive episodes \ct{Sbarra et al., 2014}
    \end{itemize}
    \item Mental health: 3$\times$ risk of psychiatric inpatient care \ct{Björkenstam et al., 2013}; higher mental-health spending \ct{Meulman et al., 2023}
    \item Physical health is mixed: more MSK and cardiovascular complaints, not significant \ct{meta-analysis of 94 studies, Pellón-Elexpuru et al., 2024}
  \end{itemize}
\end{bibframe}

\subsection{Divorce and musculoskeletal disease}
\begin{bibframe}
  \frametitle{Literature: divorce and MSK disease}
  \begin{itemize}
    \item Direct evidence is limited and contradictory
    \begin{itemize}
      \item Chronic pain is more common among the divorced, separated and widowed; MSK causes 66.8\% of chronic pain \ct{Sjøgren et al., 2009}
      \item SHARE: more chronic MSK pain among divorced women, less among men \ct{Cimas et al., 2018}
      \item HRS: entering divorce lowers the probability of high-impact pain \ct{Choi \& Wagner, 2026}; no effect in Norway \ct{Andorsen et al., 2017}
    \end{itemize}
    \item Social support is the more consistent link
    \begin{itemize}
      \item Support goes with less persistent pain; divorced or separated people without a network have more \ct{Prang et al., 2015}
      \item Lack of support: more MSK morbidity, sickness absence, failure to return to work \ct{Woods, 2005}
      \item Divorce is followed by work-limiting health, especially without remarriage \ct{Couch et al., 2015; Tamborini et al., 2016}
    \end{itemize}
  \end{itemize}
\end{bibframe}

\subsection{Musculoskeletal disease and death}
\begin{bibframe}
  \frametitle{Literature: MSK disease and death}
  \begin{itemize}
    \item MSK conditions rarely kill directly; they act indirectly
    \begin{itemize}
      \item Inflammation raises cardiovascular risk \ct{England et al., 2018}; cardiovascular mortality is about 60\% higher in rheumatoid arthritis \ct{Meune et al., 2009}
      \item Pain, then inactivity, then lower fitness and comorbidity: hip and knee osteoarthritis go with higher all-cause mortality \ct{Nüesch et al., 2011; Hawker et al., 2014}
    \end{itemize}
    \item Sickness absence due to MSK diagnoses: 1.3--1.5$\times$ all-cause mortality \ct{Jansson et al., 2012}
    \item Rehabilitation can help: work-related medical rehabilitation raises return to work by 5.8 pp \ct{Bethge et al., 2019}
    \item[$\rightarrow$] A first DRV MSK rehabilitation marks \hl{serious, work-relevant impairment} and the receipt of \hl{potentially beneficial treatment}
  \end{itemize}
\end{bibframe}

% =====================================================================
\section{Data and methodology}

\subsection{Data}
\begin{bibframe}
  \frametitle{Data: linked DRV registers, 2011--2023}
  \begin{itemize}
    \item DRV Research Data Centre: datasets linkable through a unique identifier from reporting year 2011
    \item 20\% sample: \hl{17,559,407 individuals}, \hl{185,518,738 person-years}
  \end{itemize}
  \vspace{0.2cm}
  {%
  \begin{tabularx}{\linewidth}{@{}l Y l@{}}
    \toprule
    Dataset & Content & Role \\
    \midrule
    AKVS & Actively insured persons (mostly employed) & Life course \\
    LTVS & Latently insured: contributions paused, not retired (unemployed, in education, parental leave, \dots) & Life course \\
    RTBN & Pension portfolio: pension recipients & Life course \\
    \midrule
    TD & Pension terminations due to death & Mortality \\
    RSD & Completed rehabilitations: diagnosis group, benefit type, start and end dates & MSK rehabilitation \\
    VA & Pension equalisation statistics: date of divorce, supplements and deductions & Divorce \\
    \bottomrule
  \end{tabularx}
  }
\end{bibframe}

\begin{bibframe}
  \frametitle{Merging the datasets}
  \centering
  \begin{tikzpicture}[
      font=\small,
      src/.style={draw=BiBBlau, fill=BiBBlau20, rounded corners=3pt, align=center,
                  text width=4.2cm, minimum height=3.5cm, inner sep=4pt},
      sup/.style={draw=BiBTuerkis, fill=BiBTuerkis20, rounded corners=3pt, align=center,
                  text width=4.2cm, minimum height=3.5cm, inner sep=4pt},
      arr/.style={-{Stealth[length=3mm]}, thick, color=BiBBlauK20}
    ]
    % Row 1: source datasets
    \node[src] (akvs) at (2.4,9.6)  {\textbf{AKVS}\\ actively insured};
    \node[src] (ltvs) at (7.3,9.6)  {\textbf{LTVS}\\ latently insured, not retired};
    \node[src] (rtbn) at (12.2,9.6) {\textbf{RTBN}\\ pension recipients};
    \node[sup] (td) at (17.4,9.6) {\textbf{TD}\\ deaths\\[2pt] inconsistent with BASE: 102 people dropped};
    \node[sup] (rsd)  at (22.3,9.6) {\textbf{RSD}\\ rehabilitation\\[2pt] successful, completed spells; drop $\geq 3$ spells per year ($\approx$5\%)};
    \node[sup] (va)   at (27.2,9.6) {\textbf{VA}\\ divorces\\[2pt] final divorces with pension equalisation};
    % Row 2: BASE
    \node[draw=BiBBlau, fill=BiBBlau40, rounded corners=3pt, align=center, text width=13.6cm,
          minimum height=2.0cm] (base) at (7.3,5.6)
          {\textbf{BASE} (merged by the DRV): full life course: working, between jobs, retired\\
           duplicate rows: 11.3 million excess rows dropped, most complete row kept};
    % Row 3: analysis panel
    \node[draw=BiBViolett, fill=BiBViolett20, rounded corners=3pt, align=center, text width=28.4cm,
          minimum height=1.6cm, font=\small] (panel) at (14.8,2.3)
          {\textbf{Person-year panel 2011--2023}: 17.6 million people, 185.5 million person-years};
    \foreach \n in {akvs,ltvs,rtbn} \draw[arr] (\n.south) -- (\n.south |- base.north);
    \draw[arr] (base.south) -- (base.south |- panel.north);
    \foreach \n in {td,rsd,va} \draw[arr] (\n.south) -- (\n.south |- panel.north);
    % Note
    \node[align=left, text width=28.4cm, font=\footnotesize, text=gray] at (14.8,0.6)
      {The VA misses marriages of three years or less, waived equalisation and near-equal incomes: recorded divorces fall
       from 23,364 (2011) to 9,350 (2023), faster than the national divorce rate.};
  \end{tikzpicture}
\end{bibframe}

\subsection{Matching strategy}
\begin{bibframe}
  \frametitle{Design: annual risk sets around divorce}
  \centering
  \begin{tikzpicture}[font=\footnotesize, x=1cm, y=1cm]
    \def\u{2.35}
    % positions: t0-3 at 1.2 ... t0+5 at 1.2+8u ; t0+10 shown after a break
    \foreach \k/\lab in {0/{t_0-3},1/{t_0-2},2/{t_0-1},3/{t_0},4/{t_0+1},5/{t_0+2},6/{t_0+3},7/{t_0+4},8/{t_0+5}} {
      \draw[thick] ({1.2+\k*\u},0.55) -- ({1.2+\k*\u},0.85);
      \node[below] at ({1.2+\k*\u},0.5) {$\lab$};
    }
    \draw[thick] (0.2,0.7) -- (21.6,0.7);
    \draw[thick, dashed] (21.6,0.7) -- (25.4,0.7);
    \draw[thick, -{Stealth[length=3mm]}] (25.4,0.7) -- (29.4,0.7);
    \draw[thick] (27.6,0.55) -- (27.6,0.85);
    \node[below] at (27.6,0.5) {$t_0+10$ or 2023};
    % matching window
    \fill[BiBBlau20, draw=BiBBlau, rounded corners=2pt] (0.3,2.9) rectangle ({1.2+2*\u+0.9},4.5);
    \node[align=center, font=\small] at ({1.2+\u},3.7) {\textbf{Matching}\\$t_0-3$ to $t_0-1$};
    % divorce
    \fill[BiBRot] ({1.2+3*\u},0.7) circle (0.22);
    \node[text=BiBRotK20, font=\small] at ({1.2+3*\u},1.45) {\textbf{divorce}};
    % rehab follow-up
    \fill[BiBTuerkis20, draw=BiBTuerkis, rounded corners=2pt] ({1.2+4*\u-0.9},2.9) rectangle ({1.2+8*\u+0.9},4.5);
    \node[align=center, font=\small] at ({1.2+6*\u},3.7) {\textbf{RQ1}: first MSK rehabilitation\\5 years, death as competing event};
    % mortality follow-up
    \fill[BiBViolett20, draw=BiBViolett, rounded corners=2pt] ({1.2+4*\u-0.9},1.3) rectangle (28.8,2.5);
    \node[font=\small] at (20.0,1.9) {\textbf{RQ2}: all-cause mortality, up to $t_0+10$ or 2023};
  \end{tikzpicture}

  \vspace{0.4cm}
  \raggedright
  \begin{itemize}
    \item \hl{Treated}: first observed divorce from the earliest observed marriage, cohorts 2014--2018
    \item \hl{Controls}: not yet divorced by $t_0$ (they may divorce later); no recorded end of a first marriage in or before $t_0$
    \item \hl{Both}: alive at $t_0$, records in $t_0-3$ to $t_0-1$, no MSK rehabilitation up to $t_0$
    \item Marital status is deliberately \hl{not} used: it is partly built from the VA, so ``married'' controls are disproportionately future divorcees (32\% divorced by 2018 in an earlier run)
  \end{itemize}
\end{bibframe}

\begin{bibframe}
  \frametitle{Propensity score matching}
  \compactlists
  \begin{columns}[T]
    \begin{column}{14.4cm}
      \textbf{Exact matching}: sex $\times$ 5-year age band $\times$ any rehabilitation by $t_0-1$

      \vspace{0.3cm}
      \textbf{Logit propensity score}, separately for each divorce year:
      {\small
      \begin{itemize}
        \item Age, earnings, pension income, income source, contribution periods
        \item Rehabilitation history: cumulative (mental-health) rehabilitations, years since last, starts, unsuccessful applications
        \item Days in regular and marginal employment, unemployment, care leave
        \item Children, federal state, occupation, education
        \item Missingness treated as informative (indicators)
      \end{itemize}
      }
    \end{column}
    \begin{column}{14.4cm}
      \textbf{Matching algorithm}
      \begin{itemize}
        \item Optimal 1:1 matching without replacement, caliper 0.01
        \item Divorce years in chronological order; used controls leave the pool
        \item Millions of controls: fit on all treated plus a seeded control sample (up to 500,000), correct the log-odds by $+\log f$ \ct{King \& Zeng, 2001}, then score everyone
      \end{itemize}
      \vspace{0.2cm}
      \textbf{Balance}
      \begin{itemize}
        \item Target $|\text{SMD}| < 0.10$ \ct{Austin, 2009}; any remaining imbalance would enter the outcome model
      \end{itemize}
    \end{column}
  \end{columns}
\end{bibframe}

\begin{bibframe}
  \frametitle{Why match on three pre-divorce years?}
  \begin{itemize}
    \item Matching on $t_0-1$ alone makes pairs alike in that year only; they may have arrived there along different paths
    \item Event study on lag-1 matched pairs
    \begin{itemize}
      \item Within-pair differences from $t_0-3$ to $t_0+5$, relative to $t_0-1$, standard errors clustered by pair
      \item Joint Wald test on $t_0-3$ and $t_0-2$: significant \hl{pre-trends} in earnings, rehabilitation starts and unsuccessful applications
    \end{itemize}
    \item[$\rightarrow$] Match on $t_0-3$, $t_0-2$ \textit{and} $t_0-1$: earnings, pension income, income source, unsuccessful applications, rehabilitation starts and employment days in all three years
    \item Cost: three observed pre-divorce years are needed, so only divorces in 2014--2018
  \end{itemize}
  \vspace{0.3cm}
  \begin{exampleblock}{Specifications}
    \textbf{Lag-3, 2014--2018 (preferred)} \quad$\cdot$\quad lag-1, 2014--2018 \quad$\cdot$\quad lag-1, 2012--2018
  \end{exampleblock}
\end{bibframe}

\subsection{Identification}
\begin{bibframe}
  \frametitle{RQ1: hazard of a first MSK rehabilitation}
  Follow each matched person from $t_0+1$ to $t_0+5$:
  \begin{align*}
    \lambda^{\text{rehab}}_{i,\tau} &= \Pr\left(R_{i,\tau}=1 \mid \text{no MSK rehabilitation before } \tau,\ \text{alive at } \tau\right)\\
    &= \Lam\left(\alpha_\tau + \beta D_i + \mathbf{x}_i'\boldsymbol{\gamma}\right), \qquad \tau = 1,\dots,5
  \end{align*}
  {\small
  \begin{itemize}
    \item $R_{i,\tau}$: first MSK rehabilitation in year $t_0+\tau$; $D_i$: divorce in $t_0$; $\alpha_\tau$: follow-up-year effects
    \item Three specifications for $\mathbf{x}_i$: (a) none, matching only; (b) covariates with $|\text{SMD}|>0.10$ after matching (none, so identical to (a)); (c) all covariates in $t_0+\tau-1$ (may themselves be affected by divorce)
    \item Separately for men and women; standard errors clustered by matched pair
    \item Absolute scale: 5-year \hl{cumulative incidence} with death as a competing event (discrete Aalen--Johansen); risk difference per 1,000 and risk ratio; 200 matched-pair bootstraps
  \end{itemize}
  }
\end{bibframe}

\begin{bibframe}
  \frametitle{RQ2: hazard of death}
  Follow each person from $t_0+1$ to $t_0+10$ or 2023, whichever comes first:
  \begin{align*}
    \lambda^{\text{death}}_{i,\tau} &= \Pr\left(Y_{i,\tau}=1 \mid \text{alive at the start of follow-up year } \tau\right)\\
    &= \Lam\left(\alpha_\tau + \beta D_i + \mathbf{x}_i'\boldsymbol{\gamma}\right)
  \end{align*}
  \begin{itemize}
    \item $Y_{i,\tau}$: death in year $t_0+\tau$; same specifications as RQ1, by sex
    \item Absolute scale: 5-year cumulative mortality $=1-$ Kaplan--Meier survival; risk difference per 1,000 and risk ratio
    \item Landmark analysis: condition on surviving follow-up year 1 or 2
    \item[$\rightarrow$] \hl{Mediation}: how much of the 5-year mortality gap runs through a first MSK rehabilitation?
  \end{itemize}
\end{bibframe}

\begin{bibframe}
  \frametitle{Mediation: why g-computation?}
  {\small
  \begin{itemize}
    \item A mediator measured once after divorce misses deterioration later in follow-up
    \item The post-divorce history $\mathbf{L}_{i,t}$ (earnings, pension income, income source, unsuccessful applications, non-MSK rehabilitation, employment days) is affected by divorce \textit{and} affects rehabilitation and death, so observed values cannot be plugged in
    \item[$\rightarrow$] Model the process year by year and simulate counterfactual histories \ct{Robins, 1986; Lin et al., 2017}
  \end{itemize}
  }
  \begin{block}{Yearly models, weighted by the inverse probability of having a record ($\hat\omega \leq 20$)}
    \vspace{-0.3cm}
    {\small
    \begin{align*}
      \text{Rehabilitation:}\quad \lambda^{\text{rehab}}_{i,t} &= \Lam\left(\kappa_t + \delta D_i + \mathbf{z}_i'\boldsymbol{\eta} + \mathbf{L}_{i,t-1}'\boldsymbol{\eta}_L\right)\\
      \text{Death:}\quad \lambda^{\text{death}}_{i,t} &= \Lam\left(\alpha_t + \theta_D D_i + \theta_M M_{i,t} + \theta_{DM} D_i M_{i,t} + \mathbf{z}_i'\boldsymbol{\rho} + \mathbf{L}_{i,t-1}'\boldsymbol{\rho}_L\right)\\
      \text{History:}\quad \mathbf{L}_{i,t} &\sim f\left(D_i, \mathbf{z}_i, \mathbf{L}_{i,t-1}, M_{i,t-1}, M_{i,t}\right) \quad \text{(nine models)}
    \end{align*}
    }
    \vspace{-0.3cm}
    {\footnotesize Unknown rehabilitation status in a missing year of death: average over both possibilities (50/50).}
  \end{block}
\end{bibframe}

\begin{bibframe}
  \frametitle{G-computation: simulating histories}
  \begin{columns}[T]
    \begin{column}{15.6cm}
      \centering
      \begin{tikzpicture}[
          font=\small,
          stp/.style={draw=BiBViolett, fill=BiBViolett20, rounded corners=4pt, align=left,
                       text width=12.4cm, minimum height=1.55cm, inner sep=6pt},
          arr/.style={-{Stealth[length=3mm]}, thick, color=BiBViolettK20}
        ]
        \node[draw=BiBBlau, fill=BiBBlau20, rounded corners=4pt, align=left, text width=12.4cm,
              inner sep=6pt] (init) at (0,10.2)
              {\textbf{Start}: observed history $\mathbf{L}_{i,t_0-1}$, $M_{i,t_0}=0$};
        \node[stp] (s1) at (0,7.9) {\textbf{1. Rehabilitation}: $\lambda^{\text{rehab}}_{i,t}(\mathbf{L}_{i,t-1})$, draw $M_{i,t}$};
        \node[stp] (s2) at (0,5.6) {\textbf{2. Death}: $\lambda^{\text{death}}_{i,t}(\mathbf{L}_{i,t-1}, M_{i,t})$, update survival};
        \node[stp] (s3) at (0,3.3) {\textbf{3. History}: draw $\mathbf{L}_{i,t}$ given $D_i,\mathbf{z}_i,\mathbf{L}_{i,t-1},M_{i,t-1},M_{i,t}$};
        \node[align=center, font=\footnotesize] (end) at (0,1.3) {\textbf{4.} hand $\mathbf{L}_{i,t}$ forward, $t \to t+1$, until $t_0+5$};
        \draw[arr] (init) -- (s1);
        \draw[arr] (s1) -- (s2);
        \draw[arr] (s2) -- (s3);
        \draw[arr] (s3) -- (end);
        \draw[arr] (end.west) -| ([xshift=-0.7cm]s1.west) -- (s1.west);
      \end{tikzpicture}
    \end{column}
    \begin{column}{13.2cm}
      \begin{itemize}
        \item Simulated for every matched divorcee under $d=1$ and $d=0$
        \item 25 Monte Carlo repeats per person
        \item 100 matched-pair bootstraps for confidence intervals
        \item Check first: does the simulation reproduce observed rehabilitation incidence, mortality and earnings?
      \end{itemize}
    \end{column}
  \end{columns}
\end{bibframe}

\begin{bibframe}
  \frametitle{Interventional direct and indirect effects}
  \begin{itemize}
    \item Simulated rehabilitation timing $M_i(d)$ is the output of a stochastic process, so there is no individual counterfactual for the cross-world combination $(d, M(d'))$
    \item Instead: person $i$ receives the simulated timing of a random person $j$ with the same $d$, within sex $\times$ age band $\times$ prior rehabilitation (overlap)
    \item[$\rightarrow$] Interventional effects \ct{VanderWeele et al., 2014}, not natural effects \ct{Pearl, 2001}
  \end{itemize}
  \begin{block}{Decomposition of 5-year mortality}
    \vspace{-0.4cm}
    \begin{align*}
      \text{Total effect} &= \hat Y(1,M(1)) - \hat Y(0,M(0))\\
      \text{Direct effect} &= \hat Y(1,M(0)) - \hat Y(0,M(0))\\
      \text{Indirect effect} &= \hat Y(1,M(1)) - \hat Y(1,M(0))
    \end{align*}
    \vspace{-0.4cm}
    {\small $\hat Y(d,M(d'))$: mortality with divorce status $d$ and rehabilitation timing drawn from the distribution under $d'$}
  \end{block}
\end{bibframe}

% =====================================================================
\section{Results}

\subsection{Descriptives and matching}
\begin{bibframe}
  \frametitle{Matched sample and covariate balance}
  \begin{columns}[T]
    \begin{column}{8.6cm}
      {\small
      \textbf{Sample, lag-3}\\[0.2cm]
      95,693 first divorces 2014--18\\
      $-$1,116 not first marriage\\
      $-$1,823 died or MSK rehab in $t_0$\\
      $-$576 no record in $t_0-1$\\
      $-$7,071 MSK rehab before $t_0$\\
      $-$646 no record in $t_0-2$, $t_0-3$\\[0.2cm]
      \hl{84,461 pairs} (100\% matched)\\
      41,430 men, 43,031 women\\[0.2cm]
      Mean score gap $1.1\times10^{-6}$; maximum 0.009 $<$ caliper
      }
    \end{column}
    \begin{column}{20.4cm}
      {\small
      \setlength{\tabcolsep}{4pt}
      \begin{tabular}{@{}lrrrrrr@{}}
        \toprule
        & \multicolumn{3}{c}{Before matching} & \multicolumn{3}{c}{After matching} \\
        \cmidrule(lr){2-4}\cmidrule(l){5-7}
        In $t_0-1$ unless stated & Div. & Contr. & SMD & Div. & Contr. & SMD \\
        \midrule
        Age (years) & 45.96 & 50.97 & $-0.340$ & 45.96 & 45.97 & $-0.001$ \\
        Earnings (€) & 22,511 & 14,107 & 0.412 & 22,511 & 23,709 & $-0.057$ \\
        Earnings, $t_0-3$ (€) & 20,679 & 13,119 & 0.388 & 20,679 & 22,036 & $-0.067$ \\
        Earnings missing (\%) & 13.25 & 42.76 & $-0.696$ & 13.25 & 12.62 & 0.019 \\
        Pension income (€) & 70.56 & 230.56 & $-0.438$ & 70.56 & 63.50 & 0.027 \\
        Cumulative rehabilitations & 0.151 & 0.035 & 0.320 & 0.151 & 0.150 & 0.003 \\
        Ever mental-health rehab (\%) & 5.31 & 0.78 & 0.266 & 5.31 & 4.76 & 0.025 \\
        Unsuccessful application (\%) & 2.64 & 0.43 & 0.180 & 2.64 & 2.24 & 0.026 \\
        Years since last rehab & 1.59 & 1.99 & $-0.244$ & 1.59 & 1.73 & $-0.087$ \\
        Days unemployed & 47.16 & 34.74 & 0.121 & 47.16 & 48.04 & $-0.008$ \\
        Income source (max.\ level) & & & 0.663 & & & 0.035 \\
        Occupation (max.\ level) & & & 0.448 & & & 0.024 \\
        \bottomrule
      \end{tabular}
      }
    \end{column}
  \end{columns}
  \slidenote{Mean $|$SMD$|$ over 73 terms: 0.239 before, 0.022 after matching; no covariate above 0.10 after matching.}
\end{bibframe}

\begin{bibframe}
  \frametitle{Event study after lag-3 matching}
  \compactlists
  \begin{columns}[T]
    \begin{column}{14.2cm}
      \includegraphics[width=\linewidth]{event_study_differences_lag3.png}
    \end{column}
    \begin{column}{14.8cm}
      \begin{itemize}
        \item \textbf{Before divorce}: small remaining gaps; the Wald test still rejects for earnings, rehabilitation starts and marginal employment ($p<0.001$)
        \begin{itemize}
          \item These years are matching variables, so this checks fit rather than testing independently
          \item Earnings gaps are below 1\% of mean earnings
        \end{itemize}
        \item \textbf{After divorce}, relative to $t_0-1$:
        \begin{itemize}
          \item Rehabilitation starts: +1.6 per 100 in $t_0$, +2.9 by $t_0+4$
          \item Unemployment: +11 to +15 days per year
          \item Unsuccessful applications: about +0.8 pp
          \item Marginal employment: $-5$ to $-8$ days per year
        \end{itemize}
        \item \textbf{Women}: 21.8\% had some unemployment in $t_0-1$ vs.\ 16.7\% of controls, but similar shares in $t_0-3$ and $t_0-2$: labour-market changes start before the divorce is recorded
      \end{itemize}
    \end{column}
  \end{columns}
\end{bibframe}

\subsection{RQ1: musculoskeletal rehabilitation}
\begin{bibframe}
  \frametitle{RQ1: divorce doubles the MSK rehab risk}
  \begin{columns}[T]
    \begin{column}{16.8cm}
      \includegraphics[width=\linewidth]{ms_by_sex_lag3.png}
    \end{column}
    \begin{column}{12.2cm}
      {\small
      \setlength{\tabcolsep}{4pt}
      \begin{tabular}{@{}lrr@{}}
        \toprule
        Lag-3, 2014--2018 & Men & Women \\
        \midrule
        Odds ratio, matching only & 2.04$^{***}$ & 2.30$^{***}$ \\
        Odds ratio, $t-1$ controls & 1.99$^{***}$ & 2.20$^{***}$ \\
        \addlinespace
        5-year incidence, divorced & 8.30\% & 8.81\% \\
        5-year incidence, controls & 4.16\% & 3.93\% \\
        \addlinespace
        Risk difference per 1,000 & \textbf{41.4} & \textbf{48.7} \\
        & \ci{39.0}{44.5} & \ci{46.0}{51.9} \\
        Risk ratio & 2.00 & 2.24 \\
        \bottomrule
      \end{tabular}
      }
    \end{column}
  \end{columns}
  \vspace{0.3cm}
  \begin{itemize}
    \item Pooled: 8.56\% vs.\ 4.04\%, i.e.\ \hl{45.1 more first MSK rehabilitations per 1,000}; women minus men: 7.3 per 1,000 \ci{2.8}{12.2}
    \item Robust: lag-1 specifications give odds ratios of 1.97--2.08 (men) and 2.20--2.30 (women)
  \end{itemize}
  \slidenote{$^{***}p<0.001$. Odds-ratio intervals from pair-clustered standard errors; risk intervals from 200 matched-pair bootstraps; death is a competing event.}
\end{bibframe}

\begin{bibframe}
  \frametitle{RQ1: timing of first rehabilitation starts}
  \begin{columns}[T]
    \begin{column}{16.4cm}
      \includegraphics[width=\linewidth]{annual_first_msk_starts_and_difference_by_sex_lag3.png}
    \end{column}
    \begin{column}{12.6cm}
      {\small
      \begin{itemize}
        \item \textbf{Men}: about 8--9 extra first starts per 1,000 already in year 1, stable over five years (trend $-0.08$ per year, n.s.)
        \item \textbf{Women}: the gap grows to about 11--12 per 1,000 in years 4--5 (trend $+0.78$ per year, $p<0.05$)
        \item First starts only, so repeat episodes cannot explain the widening
        \item Consistent with accumulating physical strain (no help with physical tasks after divorce)
        \item But: the year-1 gap may reflect problems that began before divorce
      \end{itemize}
      }
    \end{column}
  \end{columns}
\end{bibframe}

\begin{bibframe}
  \frametitle{RQ1: who crosses the threshold?}
  \begin{columns}[T]
    \begin{column}{15.6cm}
      \begin{block}{Risk difference per 1,000, by income source}
        {\small
        \setlength{\tabcolsep}{4pt}
        \begin{tabular}{@{}lrr@{}}
          & Men & Women \\
          \midrule
          Earnings only & \textbf{49.5} \ci{45.8}{53.1} & \textbf{55.2} \ci{51.6}{58.6} \\
          Neither observed & 10.5 \ci{2.1}{18.1} & 23.5 \ci{14.4}{32.9} \\
          Earnings and pension & 2.7 \ci{-4.3}{9.7} & 7.8 \ci{-3.9}{17.8} \\
          Pension only & $-0.1$ \ci{-1.8}{1.8} & 1.1 \ci{-1.8}{4.3} \\
        \end{tabular}
        }

        \vspace{0.1cm}
        {\footnotesize Income source in $t_0-1$. Exploratory, not yet in the draft; balance within subgroups is not guaranteed.}
      \end{block}
    \end{column}
    \begin{column}{13.4cm}
      \begin{exampleblock}{Associations with entry}
        {\small
        \begin{tabular}{@{}lr@{}}
          Unsuccessful application in $t-1$ & OR 4.23 \\
          Earnings (+1 SD, log) & OR 1.39 \\
          Pension income (+1 SD, log) & OR 0.57 \\
          Age (+1 year within band) & OR 1.03 \\
        \end{tabular}

        \vspace{0.1cm}
        {\footnotesize Model with $t-1$ controls; all $p<0.001$.}
        }
      \end{exampleblock}
    \end{column}
  \end{columns}
  \vspace{0.3cm}
  \begin{itemize}
    \item The excess is concentrated among \hl{people with earnings}, in line with DRV rehabilitation as protection of working capacity
    \item Preliminary: non-MSK rehabilitations also roughly double (mental health $\approx$2.7$\times$, nervous system $\approx$3.0$\times$), which points to broader distress
  \end{itemize}
\end{bibframe}

\subsection{RQ2: mortality}
\begin{bibframe}
  \frametitle{RQ2: excess mortality, mainly among men}
  \centering
  {\small
  \setlength{\tabcolsep}{10pt}
  \begin{tabular}{@{}lrr@{}}
    \toprule
    Lag-3, 2014--2018 & Men & Women \\
    \midrule
    Odds ratio, matching only & 1.28$^{***}$ \ci{1.22}{1.36} & 1.17$^{***}$ \ci{1.08}{1.27} \\
    Odds ratio, $t-1$ controls & 1.21$^{***}$ \ci{1.14}{1.29} & 1.05 \ci{0.95}{1.16} \\
    \addlinespace
    5-year mortality, divorced & 3.49\% & 1.25\% \\
    5-year mortality, controls & 3.12\% & 1.33\% \\
    Risk difference per 1,000 & \textbf{3.6} \ci{1.6}{5.9} & $\mathbf{-0.8}$ \ci{-2.2}{0.6} \\
    Risk ratio & 1.12 & 0.94 \\
    \addlinespace
    Deaths (full follow-up) & 4,933 & 2,050 \\
    \bottomrule
  \end{tabular}
  }\par
  \vspace{0.3cm}
  \raggedright
  \begin{itemize}
    \item Women minus men: $-4.47$ deaths per 1,000 \ci{-6.95}{-1.96}
    \item Lag-1 specifications give very similar odds ratios (men 1.25--1.28, women 1.14--1.20)
    \item The odds ratio pools up to nine follow-up years, the risk difference only five, so the two can differ in sign
  \end{itemize}
\end{bibframe}

\begin{bibframe}
  \frametitle{RQ2: early deficit, later excess}
  \begin{columns}[T]
    \begin{column}{14.4cm}
      \includegraphics[width=\linewidth]{mortality_by_sex_with_differences_lag3.png}
    \end{column}
    \begin{column}{14.6cm}
      {\small
      \begin{itemize}
        \item \textbf{Year 1}: 1.49 vs.\ 3.58 deaths per 1,000 (divorced vs.\ controls), gap $-2.1$ ($p<0.001$); men $-3.2$, women $-1.0$
        \item \textbf{Year 2}: gap $-0.65$ ($p<0.05$)
        \item \textbf{Years 3--5}: divorced men die more often: $+1.2$, $+3.6$, $+2.4$ per 1,000 (years 4--5: $p<0.001$); women n.s.
        \item The cumulative excess appears by year 4 for men and year 5 pooled; not within five years for women
      \end{itemize}
      }
    \end{column}
  \end{columns}
\end{bibframe}

\begin{bibframe}
  \frametitle{RQ2: guaranteed survival and landmarks}
  \begin{columns}[T]
    \begin{column}{15.0cm}
      {\small
      \begin{itemize}
        \item The VA records only \hl{final} divorces: a divorcee must survive until the court decision
        \item Decision after $t_0$: same year 16.5\%, one year later 64.1\%, two years 13.8\%, three or more 5.6\%
        \item People who fall seriously ill may postpone divorce or die during proceedings, and then count as controls
        \item[$\rightarrow$] Landmark analysis removes this guaranteed-survival period
      \end{itemize}
      }
    \end{column}
    \begin{column}{14.0cm}
      \begin{block}{5-year mortality risk difference per 1,000}
        {\small
        \setlength{\tabcolsep}{5pt}
        \begin{tabular}{@{}lrrr@{}}
          & From $t_0$ & \shortstack[r]{Alive end\\of year 1} & \shortstack[r]{Alive end\\of year 2} \\
          \midrule
          Men & 3.6 & 6.8 & \textbf{7.2} \\
          Women & $-0.8$ & 0.1 & 1.1 \\
          Pooled & 1.4 & 3.4 & 4.1 \\
          \addlinespace
          Risk ratio, men & 1.12 & 1.26 & 1.35 \\
        \end{tabular}
        }

        \vspace{0.1cm}
        {\footnotesize Confidence intervals still to be bootstrapped.}
      \end{block}
    \end{column}
  \end{columns}
\end{bibframe}

\begin{bibframe}
  \frametitle{RQ2: the survival gap keeps widening}
  \centering
  \includegraphics[width=24.5cm]{mortality_kaplan_meier_lag3_overall_by_sex.png}

  \raggedright
  \begin{itemize}
    \item Pooled survival at year 9: 93.07\% (divorced) vs.\ 94.91\% (controls), i.e.\ \hl{18.4 more deaths per 1,000}
    \item Years 6--9 rest on the earlier divorce cohorts only
  \end{itemize}
\end{bibframe}

\subsection{Mediation}
\begin{bibframe}
  \frametitle{Mediation: no indirect effect via MSK rehab}
  \begin{columns}[T]
    \begin{column}{15.6cm}
      {\small
      \setlength{\tabcolsep}{5pt}
      \begin{tabular}{@{}lr@{}}
        \toprule
        5-year mortality, per 1,000 (lag-3) & \\
        \midrule
        Total effect & 1.31 \ci{0.29}{1.98} \\
        Direct effect & 1.32 \ci{0.39}{2.00} \\
        Indirect effect via first MSK rehab & $\mathbf{-0.02}$ \ci{-0.16}{0.08} \\
        Proportion mediated & $-1.2$\% \ci{-20.2}{9.0} \\
        \addlinespace
        Observed difference (Kaplan--Meier) & 1.36 \ci{-0.24}{2.67} \\
        Carry-forward sensitivity: indirect & $-0.004$ \\
        \midrule
        \multicolumn{2}{@{}l}{5-year first MSK rehab, simulated vs.\ observed} \\
        Controls, $M(0)$ & 4.05\% vs.\ 4.04\% \\
        Divorced, $M(1)$ & 8.42\% vs.\ 8.56\% \\
        \bottomrule
      \end{tabular}
      }
    \end{column}
    \begin{column}{13.4cm}
      {\small
      \begin{itemize}
        \item The rehabilitation model reproduces the mediator well, so the near-zero indirect effect is not an artefact of that model
        \item Lag-1 specifications give the same picture (indirect effect $-0.02$)
        \item People who enter MSK rehabilitation die \textit{less} often within five years: about 0.9\% vs.\ 2.5\% among divorcees, 0.7\% vs.\ 2.3\% among controls
      \end{itemize}
      }
    \end{column}
  \end{columns}
  \slidenote{Interventional effects; 95\% intervals from 100 matched-pair bootstraps.}
\end{bibframe}

\begin{bibframe}
  \frametitle{Mediation: why results are still provisional}
  \begin{alertblock}{The death model under-predicts mortality}
    \begin{itemize}
      \item Simulated 5-year mortality: 1.59\% (divorced), 1.47\% (controls)
      \item Observed: 2.35\% and 2.21\%, so about \hl{7.5 deaths per 1,000 (one third) are missing}
      \item Year 1 looks right only because the errors offset: divorced $+1.04$, controls $-1.04$ per 1,000
    \end{itemize}
  \end{alertblock}
  \vspace{0.3cm}
  \textbf{Next steps}
  \begin{itemize}
    \item Add the missing earnings and employment-day variables to the history models and re-calibrate
    \item Decompose separately for men and women, as for RQ1 and RQ2
    \item Check how often a borrowed rehabilitation timing is implausible for its recipient (overlap)
  \end{itemize}
\end{bibframe}

% =====================================================================
\section{Discussion and conclusion}

\begin{bibframe}
  \frametitle{Discussion}
  \begin{itemize}
    \item Divorce roughly \hl{doubles} the 5-year risk of crossing the MSK rehabilitation threshold (+41 men, +49 women per 1,000), robust across matching specifications
    \item The excess is concentrated among people with earnings, in line with rehabilitation as protection of working capacity
    \item Women's gap widens over time, consistent with accumulating physical strain, but other rehabilitations rise too, which suggests broader distress
    \item Excess mortality among men appears once the divorce is final and keeps growing; the early deficit is plausibly a recording artefact
    \item A first MSK rehabilitation does not carry the mortality gap: entry marks treated, work-capable people rather than the frailest
  \end{itemize}
\end{bibframe}

\begin{bibframe}
  \frametitle{Mechanisms to test next}
  {\small
  \renewcommand{\arraystretch}{1.35}
  \begin{tabularx}{\linewidth}{@{}>{\raggedright\arraybackslash}p{10.4cm} Y@{}}
    \toprule
    Hypothesis & Test \\
    \midrule
    \textbf{Time to seek help}: once the marriage ends, people finally apply for rehabilitation &
      Link first starts to application dates (before or after the month the marriage ends); court decision date as a second timing marker \\
    \textbf{Loss of help with physical tasks} after separation &
      Periods of incapacity before divorce as a proxy; total rehabilitation days with and without repeat episodes \\
    \textbf{Financial urgency}: protecting one's own earning capacity &
      Work impairment at entry; people already working before the split; income quintiles \\
    \textbf{Access}: pension equalisation adds contribution periods (15-year rule) &
      Eligibility before and after equalisation \\
    \bottomrule
  \end{tabularx}
  }
\end{bibframe}

\begin{bibframe}
  \frametitle{Limitations}
  \begin{itemize}
    \item Conditional exchangeability: relationship conflict, social support and repartnering are unobserved
    \item The VA covers only final divorces with pension equalisation (no marriages of three years or less, no waived equalisation, no near-equal incomes)
    \item Controls are people who do not divorce in $t_0$: marital status is mostly unknown (83\%) and some divorce later
    \item The outcome is an administrative threshold (application and approval), not disease onset; pensioners are rarely eligible
    \item Residual pre-divorce differences in earnings, rehabilitation starts and marginal employment
    \item The mediation model is not yet calibrated and is pooled over sex
  \end{itemize}
\end{bibframe}

\begin{bibframe}
  \frametitle{Conclusion}
  \begin{block}{1. Divorce and MSK rehabilitation}
    Divorce doubles the 5-year risk of a first MSK rehabilitation: 8.6\% vs.\ 4.0\% (+45 per 1,000)
  \end{block}
  \begin{block}{2. Divorce and mortality}
    Divorce raises mortality among men; the gap grows over time (pooled: 18 more deaths per 1,000 by year 9)
  \end{block}
  \begin{block}{3. Mediation}
    So far, MSK rehabilitation does not explain the divorce--mortality gap, but the mediation model still needs re-calibration
  \end{block}
\end{bibframe}

% =====================================================================
\section*{}
\begin{bibframe}[frametype=thankyouen]
  \frametitle{Thank you!}
  \large
  \vspace{2cm}
  \textbf{Kar Man Tan}\par
  \vspace{1cm}
  Federal Institute for\\
  Population Research (BiB)\par
  \vspace{1cm}
  Phone: +49 611 75 4674\\
  karman.tan@bib.bund.de
\end{bibframe}

\end{document}
